\begin{problem}[Warmup]{(a)}
    Answer the following true or false; if your answer is false, give a counterexample. Assume that $F \subseteq K$ are fields.
    \begin{enumerate}
        \item[(a)] $\mathrm{Gal}(F \,|\, F)$ is the trivial group.
        \item[(b)] Let $\alpha \in K \setminus F$ and $\alpha$ is algebraic over $F$. Then $\mathrm{Gal}(F(\alpha) \,|\, F)$ is a finite group.
        \item[(c)] $\mathrm{Gal}(\mathbb{Q}(\sqrt{5}) \,|\, \mathbb{Q})$ is a two element group.
        \item[(d)] An automorphism preserves addition but need not preserve multiplication.
        \item[(e)] The group $\mathrm{Gal}(\mathbb{Q}(\sqrt[5]{2}) \,|\, \mathbb{Q})$ has at least five elements.
    \end{enumerate}
\end{problem}
\unsolved

\begin{problem}[Warmup]{(b)}
    Suppose that $\zeta$ is the primitive 11th root of unity. What is the Galois group $\mathrm{Gal}(\mathbb{Q}(\sqrt[11]{2}, \zeta) \,|\, \mathbb{Q}(\sqrt[11]{2}))$?
\end{problem}
\unsolved

\begin{problem}[Warmup]{(c)}
    Give examples of the following, or else argue that such an example does not exist:
    \begin{enumerate}
        \item[(a)] A Galois group with exactly 4 elements.
        \item[(b)] A field $F$ and an element $\alpha$ that is algebraic over $F$, where $\mathrm{Gal}(F(\alpha) \,|\, F)$ has infinitely many elements.
        \item[(c)] A non-cyclic Galois group.
    \end{enumerate}
\end{problem}
\unsolved

\begin{problem}{1}
    Compute the Galois group
    $\mathrm{Gal}(\mathbb{Q}(\sqrt{2}, i) \,|\, \mathbb{Q})$. Compute the Galois group
    $\mathrm{Gal}(\mathbb{Q}(\sqrt{2}, i) \,|\, \mathbb{Q}(\sqrt{2}))$.
\end{problem}
\unsolved

\begin{problem}{2}
    Consider the splitting field $\mathbb{Z}_2(\alpha)$ of the polynomial
    $x^3 + x + 1 \in \mathbb{Z}_2[x]$ that we looked at in Example 45.5. Explicitly
    describe the elements of the Galois group
    $\mathrm{Gal}(\mathbb{Z}_2(\alpha) \,|\, \mathbb{Z}_2)$ as permutations of the three
    roots $\alpha, \alpha^2, \alpha + \alpha^2$.
\end{problem}
\unsolved

\begin{problem}{3}
    Compute the Galois group $\mathrm{Gal}(K \,|\, \mathbb{Q})$, where $K$ is the
    splitting field for the polynomial $x^4 - 4x^2 + 2$. You computed this splitting field
    in Exercise 45.5. Now compute $\mathrm{Gal}(K \,|\, \mathbb{Q}(\sqrt{2}))$.
\end{problem}
\unsolved

\begin{problem}{4}
    Suppose that $\zeta$ is the primitive cube root of unity. Let $K$ be the splitting
    field of $x^9 - 1 \in \mathbb{Q}[x]$. Compute
    $\mathrm{Gal}(K \,|\, \mathbb{Q}(\zeta))$.
\end{problem}
\unsolved

\begin{problem}{5}
    Consider the polynomial $f = x^3 - 3x - 1 \in \mathbb{Q}[x]$. This is the polynomial
    we considered in Section 39.2 (and Section 44.2) when we proved that it is impossible
    to trisect a $60^\circ$ angle with ruler and compass. In this exercise we will compute
    the Galois group of the splitting field for $f$ over the rational field.
    \begin{enumerate}
        \item[(a)] Show that $f$ is irreducible over $\mathbb{Q}$.
        \item[(b)] In Section 39.2 we used a trigonometric identity to show that
        $\alpha = 2\cos 20^\circ$ is a root of this polynomial. Use similar arguments to
        show that $\beta = -2\cos 40^\circ$ and $\gamma = -2\sin 10^\circ$ are the other two
        roots of $f$.
        \item[(c)] Use elementary trigonometry to show that $\mathbb{Q}(\beta)$ is the
        splitting field for $f$ over $\mathbb{Q}$. (That is, show that
        $\alpha, \gamma \in \mathbb{Q}(\beta)$.)
        \item[(d)] We can think of $\mathrm{Gal}(\mathbb{Q}(\beta) \,|\, \mathbb{Q})$ as a
        group of permutations of the roots $\alpha, \beta, \gamma$. Argue that in this case
        this group is a cyclic group of order three.
        \item[(e)] In this problem we solved for the roots of $f$ by taking advantage of
        trigonometry. Instead, use the Cardano-Tartaglia formula to obtain a root of $f$.
        Show that the root you obtain is one of the roots given above.
    \end{enumerate}
\end{problem}
\unsolved

\begin{problem}{6}
    Apply the same reasoning as in Example 47.12 to conclude that the Galois group of the
    splitting field for $x^5 - 14x^2 + 7$ over the rational numbers is the permutation
    group $S_5$.
\end{problem}
\unsolved

\begin{problem}{7}
    Repeat the previous exercise for the polynomial
    $x^5 - 4x^4 + 2x + 2 \in \mathbb{Q}[x]$.
\end{problem}
\unsolved

\begin{problem}{8}
    Generalize the reasoning from Example 47.12 to prove the following: Consider the
    irreducible $f \in \mathbb{Q}[x]$ with $\deg(f) = p$, where $p$ is a prime integer.
    Suppose that $f$ has exactly two non-real roots. Prove that the Galois group of the
    splitting field of $f$ is the permutation group $S_p$.
\end{problem}
\unsolved

\begin{problem}{9}
    Prove the following modest variation on Theorem 47.2: Let $K$ be a simple algebraic
    extension of the field $F$ (of any characteristic). Then
    $|\mathrm{Gal}(K \,|\, F)| = [K : F]$ if and only if $K$ is a normal extension of $F$.
\end{problem}
\unsolved
